% -----------------------------------------------------------
%  Response to Reviewers — MIDS (TIFS Major Revision)
%
%  Compile with pdflatex; depends only on standard packages.
%  Edit the manuscript ID and date below before submitting.
% -----------------------------------------------------------

\documentclass[11pt,a4paper]{article}

\usepackage[a4paper, margin=2.4cm]{geometry}
\usepackage[T1]{fontenc}
\usepackage[utf8]{inputenc}
\usepackage{lmodern}
\usepackage[dvipsnames]{xcolor}
\usepackage{enumitem}
\usepackage{hyperref}
\usepackage{microtype}
\setlist{nosep}

% Visual conventions
\newcommand{\rcomment}[1]{%
  \par\medskip
  \begingroup
  \itshape
  \textbf{Reviewer comment.}\,\,#1
  \endgroup
  \par
}
\newcommand{\rresponse}[1]{%
  \par
  \begingroup
  \color{NavyBlue}
  \textbf{Response.}\,\,#1
  \endgroup
  \par\medskip
}
\newcommand{\xref}[1]{\textbf{#1}}

\title{Response to Reviewers \\
       \large MIDS: Detecting Stealthy Masquerade and Tampering
       Attacks on CAN Bus via Bidirectional Mamba}
\author{}
\date{}

\begin{document}
\maketitle
\vspace{-3em}

\noindent\textbf{Previous Manuscript ID:} \texttt{T-IFS-21023-2025} \\
\noindent\textbf{Current Manuscript ID:} \texttt{T-IFS-25937-2026} \\
\noindent\textbf{Submission:} Resubmission after Reject decision of 16-Jul-2025 \\

\hrule
\medskip

\noindent Dear Editor and Reviewers,

We thank the editor and the three reviewers for the rigorous and
constructive feedback on our previous submission. The revised
manuscript addresses every concern raised; this letter maps each
reviewer comment to the corresponding revision and indicates the
section, table, or figure of the revised paper where the change
can be inspected.

\medskip
\noindent The four most substantive revisions are:
\begin{itemize}[leftmargin=*]
\item A \textbf{new cross-dataset benchmark} (\xref{Table~VII},
\xref{Section~VI-F}) in which all eight reproduced baselines are
re-trained from scratch on four public datasets (OTIDS, ROAD,
CrySyS, CT\&T) plus our Tesla Model~3 dataset, under a single
unified protocol with 5-fold block-shuffled cross-validation.
A total of \textbf{160} runs underpin the new table.
\item A \textbf{formal threat model and attack taxonomy}
(\xref{Section~IV}) that adopts the Cho--Shin masquerade/fabrication
naming convention and explicitly describes attacker capabilities,
in-situ substitution semantics, hardware-assisted CRC bypass, and
the offline synthesis of evaluation traces.
\item Two new subsections on \textbf{data-leakage prevention and
the evaluation protocol} (\xref{Section~VI-A4} and \xref{VI-A5}),
disclosing the block-shuffled 5-fold split, the $L$-frame buffer,
the cross-scenario interleaving, and an empirical scenario-leakage
verification.
\item A new \textbf{Computational Efficiency and Deployability}
section (\xref{Section~VI-G}, \xref{Table~VI}) reporting parameters,
FLOPs, single-window latency, and peak memory for MIDS and the
strongest baselines.
\end{itemize}

In the rest of this letter, comments are quoted verbatim in
italics and our responses follow in blue. Cross-references to
sections, tables, and figures point to the \emph{revised}
manuscript.

\bigskip

% ====================================================================
\section*{Response to Reviewer 1}
\addcontentsline{toc}{section}{Response to Reviewer 1}

We thank Reviewer~1 for the supportive and detailed feedback. The
seven concerns are addressed below.

\rcomment{The paper lacks a proper description of an attacker model.
It is unclear what tampering capabilities the attackers possess
and what specific modifications in the CAN signals the IDS targets.}
\rresponse{We have added \xref{Section~IV} (``Threat Model and
Attack Taxonomy''), a dedicated section formalising the attacker
model. \xref{Section~IV-A} enumerates three attacker capabilities
(Software-level Control, In-situ Modification, Hardware-assisted
CRC bypass). \xref{Section~IV-B} catalogues the three concrete
attack vectors (ID-only masquerade, Data-only tampering, Combined
masquerade) with worked Tesla examples (e.g.,~substituting the
door-status ID \texttt{0x102} with the steering-angle ID
\texttt{0x132}). \xref{Section~IV-C} (``Implications for
Detection'') then derives the two design requirements
(frequency-invariance robustness and cross-field reasoning) that
motivate MIDS's dual-stream architecture.}

\rcomment{Recent literature commonly uses a specific naming
convention for attacks to avoid confusion. I recommend that the
authors adopt this naming convention (Cho \& Shin, 2016).}
\rresponse{The Cho--Shin convention (Cho and Shin, USENIX Security 2016) is now adopted
throughout the manuscript. Specifically, \xref{Section~IV} opens
by explicitly distinguishing \emph{fabrication} attacks (frames
appended to the bus, disturbing per-ID periodicity) from
\emph{masquerade} attacks (legitimate frame suppressed and
substituted in situ, preserving periodicity). All three attack
vectors considered in this paper are instances of the masquerade
pattern, which we make explicit in the opening paragraph of
\xref{Section~IV} and in the dataset-synthesis discussion in
\xref{Section~VI-A3}.}

\rcomment{There are significant issues with the presentation of
the attacks. [\dots] Two of the presented tampering strategies
modify the ID field, effectively turning these tampering attacks
into injection attacks. This presents a significant contradiction
to the paper's original positioning.}
\rresponse{This was a critical misunderstanding caused by our
earlier writing, which we have now corrected. Under the Cho--Shin
taxonomy adopted in the revised \xref{Section~IV}, modifying the
ID field is \emph{not} equivalent to an injection attack: an
injection attack appends an extra frame to the bus (raising the
per-ID frame rate), whereas our ID-only masquerade
\emph{substitutes} a legitimate frame at its original
transmission slot (preserving the per-ID rate). We have rewritten
the opening paragraph of \xref{Section~IV} to make this in-situ
semantics explicit, and \xref{Section~IV-B} now describes all
three masquerade variants under this consistent framing.}

\rcomment{Another serious problem is that the evaluation lacks any
comparison with prior work. [\dots] They should run their approach
on at least some of these datasets to enable a comparison of
results.}
\rresponse{This is now the central new contribution of the
revision. \xref{Section~VI-F} reports MIDS's performance on four
public datasets (OTIDS, ROAD, CrySyS, CT\&T) using
\xref{Table~VI} (per-metric MIDS breakdown). In addition,
\xref{Table~VII} (``Cross-dataset F1\,(\%) of MIDS against eight
baselines under the unified 5-fold protocol'') reports F1 for
\emph{all} eight reproduced baselines on \emph{all} five datasets,
totalling 160 runs. The baseline reproduction procedure is
documented in \xref{Section~VI-B3} (``Baseline Reproduction
Protocol''); the full code is publicly released.}

\rcomment{Since the proposed method would be deployed in a
resource-constrained embedded environment, the paper should
include a brief discussion of the method's resource requirements.}
\rresponse{We added \xref{Section~VI-G} (``Computational Efficiency
and Deployability'') with \xref{Table~VI}, reporting parameter
count, FLOPs per forward pass, single-window latency (mean of
1000 timed runs after 100 warm-up runs, batch size 1), and peak
GPU memory for MIDS and the strongest reproducible baselines on
an RTX~4090. MIDS attains $1.147$~ms per 100-frame window,
$8\times$ below the fastest periodic CAN inter-arrival in our
Tesla traces, and a peak memory of $27$~MB---both well within the
envelope of contemporary automotive gateway ECUs.}

\rcomment{In Table~2, the CrySyS CAN Dataset should not be listed
under the ``simulated'' category but under ``real \& simulated''.}
\rresponse{Corrected. \xref{Table~II} now lists CrySyS under
``Real \& Simulated'', matching the published characterisation
of the dataset.}

\rcomment{There are also a few minor typesetting errors. For
instance, in the first paragraph of the Introduction section, on
line~39, a space is missing before the sentence beginning with
``However, the protocol's\dots''.}
\rresponse{Corrected, together with a careful pass over the
manuscript for similar minor issues.}

% ====================================================================
\section*{Response to Reviewer 2}
\addcontentsline{toc}{section}{Response to Reviewer 2}

We thank Reviewer~2 for the rigorous critique. The reviewer raised
seven substantive concerns. All seven have been addressed through
the new cross-dataset benchmark (\xref{Table~VII},
\xref{Section~VI-F}), the explicit disclosure of train/test/window
parameters (\xref{Section~VI-A5}, \xref{VI-B3}), the formal threat
model (\xref{Section~IV}), and updated discussion of recent
literature including the four references the reviewer suggested.

\rcomment{The paper claims they are able to outperform other
models; however, the results shown were tested on other datasets
(other than the authors' own model), as far as I can infer
[\dots]. So these results, unless the authors recreated them using
the models on their dataset, are questionable.}
\rresponse{We have made the reproduction methodology explicit in
\xref{Section~VI-B3}. All baseline numbers in our manuscript---both
on our Tesla dataset (\xref{Table~IV}) and on the four public
datasets (\xref{Table~VII})---are obtained by re-training every
baseline from scratch on the target dataset under a single
unified protocol: 50 epochs, Adam, $\mathrm{lr}=10^{-4}$, cosine
annealing with $T_{\max}=10$, gradient clipping at $\ell_2$-norm
$1.0$, batch size 1024, 5-fold block-shuffled cross-validation,
identical input tensor shape $(B,L{=}100,F{=}9)$ for every model,
and identical four-class output head $\{$Normal, ID, Data,
Both$\}$. No baseline shares weights with MIDS or borrows numbers
from its originating publication. The complete re-implementation
is released in our public repository.}

\rcomment{There are a number of IDS papers (see [1] to [4] for
example), that has superior performance than the SOTA models
listed in the paper [\dots].}
\rresponse{We have integrated all four references into
\xref{Section~II-A}: Khandelwal and Shreejith (ASAP 2023),
Park~\textit{et al.}\ (IEEE Access 2023),
Ma~\textit{et al.}\ (Security and Communication Networks 2022), and
Rangsikunpum~\textit{et al.}\ (Journal of Systems Architecture 2024)
are now discussed as a
coherent line of work pursuing FPGA-, graph-classification-, and
lightweight-GRU-based detectors. We respectfully note that all
four are evaluated exclusively against fabrication-style
(injection) threat models, where per-ID inter-arrival statistics
provide the primary detection signal. Under the masquerade model
adopted in our work (\xref{Section~IV}) this signal vanishes by
construction, so their headline F1 figures on injection
benchmarks are not directly comparable to MIDS's F1 on
masquerade benchmarks; the revised related-work discussion makes
this distinction explicit.}

\rcomment{[\dots] the data spoofing attack in the Open CAN dataset
has the same signature as the tampering attack in this case, so
the performance could be comparable on this dataset.}
\rresponse{This observation is an empirical prediction that our
new experiments directly test. On OTIDS (Open CAN-derived
benchmark), \xref{Table~VII} shows MIDS at F1 = $99.61\%$ while
the strongest baselines (DCNN, CanTransformer) both reach the
$100\%$ ceiling, so the ranking is ambiguous at saturation. To
provide a sharper comparison, the revision evaluates four public
benchmarks; on the harder masquerade-style benchmarks ROAD and
CrySyS, MIDS exceeds the strongest baseline by $+13.94$ and
$+6.80$ percentage points respectively, providing exactly the
discriminating evidence the reviewer's hypothesis requires.}

\rcomment{The testing methodology is not clear---are the results
reported on the training set [\dots]? What is your
train-validation-test split [\dots]? Do you also use data balancing
for your test set? What is the batch size / event-window size for
your test sets?}
\rresponse{We added \xref{Section~VI-A5} (``Data Leakage Prevention
and Evaluation Protocol''), which states the protocol explicitly:
5-fold block-shuffled chronological cross-validation; in each
fold the data are split 80\% training / 20\% test along contiguous
temporal blocks, with no validation split used for tuning (all
hyperparameters fixed across folds and listed in
\xref{Table~III}). Reported metrics are computed exclusively on
the held-out 20\% test block of each fold and averaged across the
five folds. Window length is $L=100$ frames with stride $S=L=100$
(non-overlapping); batch size is 1024 for both training and
inference. Dynamic class weighting is applied only to the
training loss; the test set retains its natural label
distribution.}

\rcomment{In your dataset, how do you label an event as an
intrusion (at a message level or a block level)? [\dots] A tampered
message requires recalculation of CRC and so on to be considered
as a valid message, so your injection setup integrates a
malicious ECU?}
\rresponse{Both points are now addressed in \xref{Section~IV}.
\emph{Labelling level:} attack labels are assigned at the
window level (100 frames). \emph{Tampering mechanism and CRC:}
the third bullet of \xref{Section~IV-A} (``Hardware-assisted
CRC'') explains that the attacker has compromised the software
stack of an ECU connected to the bus; the on-chip CAN controller
of that ECU automatically computes and appends a valid CRC over
the tampered payload before transmission, so receiving nodes
accept the malicious frames as protocol-compliant. We also
clarify in a new paragraph that the attack traces used for
evaluation are synthesised \emph{offline} from benign captures,
yielding byte-exact ground-truth labels, and that real-time
dominant-bit overwriting on a live bus is a separate hardware
engineering problem that is not the subject of this detector.}

\rcomment{What is the test platform? What is your inference
performance (messages per second [\dots])? What is the energy
consumption [\dots]?}
\rresponse{\xref{Section~VI-G} reports the test platform (NVIDIA
RTX~4090, batch size 1, 1000 timed runs after 100 warm-up runs)
and \xref{Table~VI} summarises params, FLOPs, single-window
latency, and peak memory for MIDS and the strongest baselines.
The MIDS latency of $1.147$~ms per 100-frame window corresponds
to a throughput of approximately $8.7\times 10^{4}$ messages/s
under windowed inference; the slowest periodic CAN signals in
our Tesla traces have inter-arrival periods of $\sim 10$~ms,
leaving an $8\times$ deployability headroom. Energy measurement
at the gateway-ECU level requires hardware power-rail
instrumentation that is part of our planned follow-up work; we
acknowledge this as a limitation in \xref{Section~VII} (``Future
Works'').}

\rcomment{Could you show your model's performance on other
datasets (especially looking at spoofing data) to see if the
performance is generalisable?}
\rresponse{This is now the central new evaluation result of the
revision. \xref{Table~VII} reports F1 for all nine models on all
five datasets, including the spoofing-style OTIDS and CT\&T
datasets specifically requested by the reviewer. MIDS achieves
the highest F1 on every dataset, with margins ranging from a
near-tie on saturated OTIDS to $+13.94$ percentage points on
ROAD. \xref{Section~VI-F} discusses the cross-dataset pattern in
detail and observes that MIDS is the only model that ranks first
or tied-first under both dense-attack (Tesla, OTIDS, CT\&T) and
sparse-attack (ROAD, CrySyS) regimes.}

% ====================================================================
\section*{Response to Reviewer 3}
\addcontentsline{toc}{section}{Response to Reviewer 3}

We thank Reviewer~3 for the careful and constructive feedback and
for highlighting both major and minor improvements. The eight
major concerns and two minor suggestions are addressed below.

\rcomment{Concern 1. The dataset is organized into 4 folders by
driving scenario [\dots] rather than by the model's required
output classes [\dots]. This creates a high risk of data leakage
[\dots]. The authors must provide a detailed explanation of their
cross-validation strategy that explicitly demonstrates how they
prevented this form of data leakage [\dots].}
\rresponse{We have added two new subsections that jointly address
this concern. \xref{Section~VI-A4} (``Decoupling Scenario from
Attack Label'') explains the two pipeline choices that prevent
scenario--label correlation: (i) within-scenario attack
replication, where every attack strategy and intensity level is
instantiated within every driving scenario, and (ii)
cross-scenario interleaving, where per-scenario recordings are
concatenated and split into chunks of 1000 windows that are
uniformly shuffled before fold partitioning.
\xref{Section~VI-A5} (``Data Leakage Prevention and Evaluation
Protocol'') details the temporal-leakage safeguards
(non-overlapping windows with $S=L=100$, $L$-frame safety buffer
between adjacent training and testing blocks) and reports an
empirical scenario-leakage verification: a held-out classifier
trained to predict the driving scenario from a window achieves
accuracy indistinguishable from chance after our pipeline,
whereas the same classifier achieves $>97\%$ accuracy on the
unshuffled raw recordings. This confirms that scenario
information is not a detectable signal in the data presented to
MIDS.}

\rcomment{Concern 2. The paper is not reproducible in its current
form. [\dots] The authors must provide the fully labeled dataset
[\dots] and add a dedicated ``Data Preprocessing'' subsection
that details the exact transformation from raw logs to
model-ready tensors.}
\rresponse{We have made the pipeline fully reproducible.
\xref{Section~VI-A5} now describes the windowing and labelling
transformation (sliding window with $S=L=100$, window-level
majority-vote labelling, block-shuffled fold construction with
$1000$-window chunks). The full preprocessing code---including
the \texttt{.asc}-to-\texttt{.csv} conversion script, the attack
synthesis tool that produces byte-exact ground-truth labels, the
windowing module, and the fold-generation utility---is released
in our public repository together with the full labelled dataset
in CSV form (one label per window). The repository URL is given
in the footnote of \xref{Section~I-B}.}

\rcomment{Concern 3. The description of the model's architecture
is incomplete. The method for fusing the forward and backward
Mamba outputs is described only as a ``weighted integration
mechanism'' [\dots]. Provide the precise mathematical formulation
[\dots] and update the model architecture diagram (Fig.~3) to
explicitly illustrate this mechanism.}
\rresponse{We provide the precise fusion formulation in
\xref{Section~V-D}: the global representation is
$\mathbf{Y}_{\mathrm{global}} = \alpha\cdot\mathbf{H}_{\mathrm{fwd}}
+ \beta\cdot\mathbf{H}_{\mathrm{bwd}}$, where $\alpha$ and
$\beta$ are scalar learnable parameters jointly optimised with
the rest of the network. The fused representation is then
flattened and passed through a fully connected layer followed by
softmax: $\hat{y} = \mathrm{Softmax}\!\left(\mathbf{W}_{\mathrm{out}}
\cdot\mathrm{Flatten}(\mathbf{Y}_{\mathrm{global}})
+ \mathbf{b}_{\mathrm{out}}\right)$. \xref{Figure~3} (model
architecture) has been redrawn to show the two learnable scalars
$\alpha$ and $\beta$ explicitly at the merge point of the
forward and backward branches.}

\rcomment{Concern 4. The paper fails to demonstrate the
generalizability of the MIDS framework. [\dots] evaluate the MIDS
framework on at least one other public benchmark dataset
containing tampering attacks, such as the CrySyS dataset.}
\rresponse{The revision evaluates MIDS on \emph{four} public
benchmarks: ROAD, CrySyS, OTIDS, and CT\&T (\xref{Table~VI} and
\xref{Section~VI-F}). MIDS attains F1 from $93.70\%$ (CT\&T) to
$99.61\%$ (OTIDS), with $95.76\%$ on the CrySyS benchmark
specifically requested by the reviewer. The cross-dataset
performance is further contextualised in \xref{Table~VII} against
eight reproduced baselines under the unified protocol.}

\rcomment{Concern 5. The experimental comparison is incomplete.
The set of baselines is missing a crucial comparison against
modern Transformer-based models [\dots].}
\rresponse{CanTransformer (Jo and Kim, IEEE Access 2024), a recent Transformer-based
IDS, has been added to the baseline set. It now appears in every
comparison table (\xref{Table~IV} on Tesla and \xref{Table~VII}
on the cross-dataset benchmark). \xref{Section~VI-D} discusses
its performance explicitly: CanTransformer is the strongest
baseline on Tesla (F1 = $88.66\%$, second only to MIDS at
$96.94\%$) and ties for the strongest baseline on OTIDS and
CT\&T, supporting the framing of Mamba versus attention as the
key architectural comparison.}

\rcomment{Concern 6. The evaluation relies on overly simplistic
metrics for a security application. Aggregate scores like F1-score
can hide an unacceptably high False Positive Rate [\dots]. Provide
a more detailed analysis [\dots] specifically reporting the FPR
and FNR. Include a discussion on the practical trade-offs,
potentially supplemented by a ROC curve.}
\rresponse{We now report FPR and FNR for MIDS on every dataset in
\xref{Table~VI} (FPR $\le 1.03\%$ and FNR $\le 8.49\%$ across all
five evaluated benchmarks; the elevated FNR on CT\&T reflects its
$>95\%$-majority Normal class, an artefact of the original
release rather than a model limitation). We also include per-class
ROC curves and per-class AUC in \xref{Figure~6(c)}, showing that
MIDS achieves AUC $> 0.998$ in every attack class on our Tesla
dataset. \xref{Section~VI-F} discusses the FPR/FNR trade-off
explicitly.}

\rcomment{Concern 7. The claim of a ``novel'' model is an
overstatement. [\dots] Suggestion: Rephrase to use the more
accurate term ``innovative framework''.}
\rresponse{We have adopted the reviewer's suggested phrasing
throughout the manuscript. Architectural claims now use
``innovative framework''; the term ``novel'' is reserved for
contributions that are genuinely new, namely the publicly
released Tesla dataset and the first-time application of Mamba
to CAN-bus masquerade detection.}

\rcomment{Concern 8. The description of the training and testing
protocol lacks explicit detail. The paper mentions a ``5-fold
cross-validation'' approach but does not explicitly state the
train/test split ratio (e.g., 80/20) [\dots].}
\rresponse{\xref{Table~III} now lists ``Train/Test split: 4{:}1''
(i.e., 80\%/20\%) explicitly. \xref{Section~VI-A5} adds a sentence
stating: ``In each fold, four blocks (80\%) form the training set
and the remaining block (20\%) is held out for testing, so the
test traffic is always chronologically distinct from the training
traffic.''}

\medskip
\noindent\textit{Minor suggestions.}

\rcomment{Minor 1. [\dots] the models listed (Attention, CNN,
LSTM, etc.) are architectures, not specific published methods.
It would be stronger to cite the specific papers from which these
results were reproduced for clarity as some references are missing
for KNN, Attention and MLP.}
\rresponse{We have replaced the generic architecture labels in
\xref{Table~IV} with specific published method names and
citations (GIDS [Seo \textit{et al.}\ PST 2018],
CANShield [Shahriar \textit{et al.}\ IoT-J 2023],
CanBus-IDS [Hoang and Kim, Veh.\ Commun.\ 2022],
DCNN [Song \textit{et al.}\ Veh.\ Commun.\ 2020],
CANTransfer [Tariq \textit{et al.}\ SAC 2020],
CanTransformer [Jo and Kim, IEEE Access 2024]),
and folded the foundational MLP/CNN
into a separately labelled block. \xref{Table~VII} maintains the
same convention.}

\rcomment{Minor 2. The term ``spoofing'' is used in the caption
for Fig.~7, but the rest of the paper uses ``tampering''. Please
ensure consistent terminology throughout.}
\rresponse{Corrected. Figure 7's caption and all related class
labels now use ``ID tampering'', ``data tampering'', and
``combined tampering'' consistently with the rest of the
manuscript.}

% ====================================================================
\hrule
\medskip

\noindent We are grateful to all three reviewers for the depth
of their engagement with the manuscript. The cross-dataset
benchmark, the formal threat model, and the explicit reproduction
protocol introduced in this revision were directly motivated by
their comments and substantially strengthen the paper. We hope
that the changes meet the reviewers' expectations.

\bigskip

\noindent Sincerely, \\
The Authors

\end{document}
